% ============================================================
% Section 2: Literature Review
% ============================================================
\section{Literature Review}

\subsection{Overview of Existing Safety Solutions}
Digital interventions for women's safety have evolved significantly over the past decade, broadly falling into three categories as shown in Table~\ref{tab:existing_solutions}.

\begin{table}[htbp]
\caption{Categories of Existing Safety Solutions}
\label{tab:existing_solutions}
\centering
\scriptsize
\begin{tabularx}{\columnwidth}{|p{1.6cm}|p{2.2cm}|X|}
\hline
\textbf{Category} & \textbf{Examples} & \textbf{Core Mechanism} \\
\hline
Traditional SOS Apps & bSafe, Circle of 6, Nirbhaya App & Manual panic button, SMS/call to contacts \\
\hline
GPS Tracking \& Geofencing & Life360, FamiSafe, Google Trusted Contacts & Continuous location sharing + zone-based alerts \\
\hline
AI/ML-Enhanced Prototypes & SAFECITY, SHIELD, VAMA & Sensor fusion + behavioral modeling + predictive alerts \\
\hline
\end{tabularx}
\end{table}

\subsection{Critical Analysis of Limitations}

\subsubsection{Traditional SOS Applications}
While widely adopted, manual-trigger systems suffer from fundamental usability gaps in real emergencies: (a)~\textit{Interaction Dependency}---require conscious, deliberate user action, often impossible during physical restraint, panic-induced freezing, or device confiscation \cite{who2018}; (b)~\textit{Contextual Blindness}---treat all locations, times, and movement patterns identically, ignoring environmental risk factors \cite{dhillon2021}; (c)~\textit{Alert Fatigue}---high false-positive rates from accidental triggers reduce trust and response urgency among contacts \cite{uddin2022}.

\subsubsection{GPS-Centric Tracking Systems}
Location-sharing platforms improve situational awareness but introduce new constraints: \textit{Privacy Concerns}---continuous location broadcasting raises surveillance and data-misuse risks \cite{gupta2023}; \textit{Battery Drain}---high-frequency GPS polling significantly reduces device usability \cite{nguyen2023}; \textit{Reactive Logic}---geofence breaches are detected post-entry, not pre-escalation \cite{gemini2025}.

\subsubsection{AI-Based Research Prototypes}
Emerging academic work explores predictive safety using machine learning. SAFECITY aggregates anonymous harassment reports to generate community heatmaps but lacks real-time individual risk assessment \cite{chen2024}. SHIELD employs accelerometer-based fall detection but operates in isolation from contextual cues like time or location type \cite{gupta2023}. VAMA integrates voice-command SOS but depends on clear audio input, which may be unavailable in noisy or constrained scenarios \cite{park2023}.

\subsection{Identified Research Gaps}
Despite technological advances, no existing solution comprehensively addresses the following requirements simultaneously: (1)~Proactive threat detection before manual intervention; (2)~Contextual intelligence integrating temporal, spatial, and behavioral signals; (3)~Autonomous, user-independent alerting; (4)~Production-ready privacy with offline resilience; (5)~Scalable cloud-native architecture supporting real-time AI inference.

SafeHer directly addresses these gaps through its Verified Safety Matrix---a hybrid engine fusing rule-based logic with large language model (LLM) contextual reasoning---deployed on a privacy-first, battery-optimized mobile architecture.
